\section{Community Detection Methods}
In this section, we will explore the community detection methods that form
the core of this study.
\subsection{Edge Betweenness}
The edge betweenness method, introduced by Girvan and Newman (2002),
was made to sidestep the shortcomings of traditional community detection methods like 
the hierachical clustering approach. The hierachical clustering method focuses on 
trying to construct a measure that tells us which edges are most central to communities.
Using these measures, it constructs communities by adding the strongest edges
to an initiall empty graph. However, this method has several drawbacks,
such as the fact that it relies on a subjective choice of the similarity measure.
Furthermore, because it tries to build communities by adding edges,
early mistakes made by noise in the data cannot be corrected later on.
The edge betweenness method, on the other hand, removed weak edges from the 
original graph, instead of adding edges.
To find the weak edges of a graoh, Girvan and Newman generalized Freeman's betweenness centrality
to edges. \textcite{newmanandgirvan2002} states if a network contains communities
or groups that are only loosely connected by fe integroup edges, then all shortest paths
different communities must go along one of these few edges. Thus, the edges connecting
communities will have high edge betweenness. By removing these edges, we seperate groups
from one another and so reveal the underlying community structure of the graph.
The algorithm they propose is simply stated as follows:
\begin{enumerate}
  \item Compute the edge betweenness for all edges in the network.
  \item Remove the edge with the highest edge betweenness.
  \item Recalculate the edge betweenness for all edges affected by the removal.
  \item Repeat from step 2 until no edges remain.
\end{enumerate}
The iteration ends when no edges remain, so in praxis, we will have
to decide when to stop the algorithm. This can be done by looking at the modularity of
the network after each iteration and stopping when the modularity is maximized.
This algorithm is computationally expensive, because we have
to recalculate the edge betweenness after each iteration. Making the algorithm
runs in worst-case time $O(m^2n)$, where $m$ is the number of edges and $n$ is 
the number of nodes in the graph.
 So it is tempting to calculate the betweenness of all edges only once and then remove them in order
of decreasing betweenness. However, we find that this approach does not work well,
because if two communities are connected by more than one edge, then there is 
no guarantee that all of the edges will have high betweenness. So by
recalculating betweennesses after the removal of each edge, we can ensure that 
at least one of the remaining edges between two communities will always have high betweenness.

\subsection{Louvain Method}
The Louvain Algorithm 




